% Glossary terms from ch13: lines of the form \item[Term] Definition.
\item[Reciprocal velocity obstacle (RVO)] The set $\RVO_{A|B}(\vel_B,\vel_A)$ of velocities $\vel$ for which $2\vel-\vel_A$ lies in the velocity obstacle $\vel_B+\VO_{A|B}$: the velocities that $A$ must avoid when it takes half of the avoidance and trusts $B$ to take the other half; it is the velocity obstacle cone with its apex moved to $\tfrac12(\vel_A+\vel_B)$.
\item[Reciprocal dance] The oscillation of two agents that avoid each other with plain velocity obstacles (both swerve, both return, both swerve) or, for RVO, that disagree about the side on which to pass; \orca removes it by fixing the share of the responsibility.
\item[Truncated velocity obstacle] $\VO^{\ttc}_{A|B}$, the set of relative velocities that lead to a collision within the horizon $\ttc$: the velocity obstacle cone cut off by the disk of center $\pos_{\mathrm{rel}}/\ttc$ and radius $R/\ttc$; convex, with a boundary made of a front arc and two legs.
\item[ORCA half-plane] $\ORCA^{\ttc}_{A|B}=\set{\vel : (\vel-(\vel_A+\tfrac12\vect{u}))\cdot\vect{n}\ge0}$, the velocities permitted to $A$ with respect to $B$ for the horizon $\ttc$, where $\vect{u}$ is the smallest change of the relative velocity that reaches the boundary of $\VO^{\ttc}_{A|B}$ and $\vect{n}$ the outward normal there; in three dimensions a half-space.
\item[Responsibility share] The fraction $\alpha$ of the change $\vect{u}$ that an agent takes: $\tfrac12$ between agents that both run \orca, $1$ (full responsibility) against an obstacle or a non-cooperative intruder; the pairwise guarantee holds exactly when the two shares of a pair sum to one.
\item[Preferred velocity] $\vel^{\mathrm{pref}}$, the velocity an agent would fly if it were alone, typically toward the next waypoint of its nominal path at cruise speed; \orca returns the permitted velocity closest to it.
\item[Time horizon (ORCA)] $\ttc$, the time span for which the chosen velocities are guaranteed collision-free; recomputed every $\dt\ll\ttc$, with a shorter horizon $\ttc_{\mathrm{obst}}$ for static obstacles.
\item[Incremental linear program] The randomized algorithm that finds the point of an intersection of half-planes and a disk closest to a given point by adding the half-planes one at a time and, whenever the current optimum violates the new one, solving a one-dimensional program on its boundary line; expected linear time in the number of half-planes.
\item[Dense fallback] The remedy for an empty intersection of \orca half-planes: choose the velocity that minimizes the largest violation over all agent half-planes while keeping the obstacle half-planes hard; van den Berg et al.\ call it the safest possible velocity, the one which penetrates the half-planes least. It is a heuristic: it minimizes a penetration, not the time to the first collision.
\item[Non-cooperative obstacle] A neighbor that does not avoid, such as a static obstacle or an unknown intruder flying its predicted velocity; the agent takes full responsibility for it, with a radius inflated by the uncertainty of the prediction.
